\Needspace{12\baselineskip}
\fancyhead[L]{\small ERSTER KOMET VON 1805.}
\begin{center}
ERSTER KOMET DES JAHRES 1805.

\rule{0.6\linewidth}{0.4pt}

Monatliche Correspondenz. Band XIII. S. 79–83. Januar 1806.

\rule{0.6\linewidth}{0.4pt}

—\quad—\quad—
\end{center}

Dr. \textsc{Gauss} berichtet uns am 5. Dezember aus Braunschweig, er habe vergeblich versucht, diese Elemente von \textsc{Bessel} an die Beobachtungen aus Frankfurt und Marseille anzupassen. Diese Beobachtungen scheinen allerdings sehr ungenau zu sein. „Da diese Elemente“, schreibt Dr. \textsc{Gauss}, „um mehrere Grad von \textsc{Huth}s letzten Beobachtungen abweichen, versuchte ich, sie anhand dieser Beobachtungen zu verbessern, obwohl \textsc{Huth}s Angaben, die nur in Schritten von fünf Minuten angegeben sind, bloße Schätzungen nach Augenmaß zu sein scheinen. Der Versuch gelang jedoch nicht. Auch die Beobachtungen von \textsc{Thulis}, die Sie mir freundlicherweise mitgeteilt haben (Novemberheft 1805, \href{https://www.digitale-sammlungen.de/en/view/bsb10538604?page=518}{S. 502}), scheinen nicht besonders gut zu sein. Jedenfalls konnte ich auch mit ihnen kein zufriedenstellendes Ergebnis erzielen. Das zeigt schon der Vergleich mit \textsc{Olbers}’ Beobachtungen. Rechnet man sie auf dieselbe Zeit um, so weichen \textsc{Olbers}’ Beobachtungen am 29. Oktober um sechs Minuten in der Rektaszension und neun Minuten in der Deklination, am 31. Oktober um eine Minute in der Rektaszension und zwölf Minuten in der Deklination von denen von \textsc{Thulis} ab. Am 9. November weichen \textsc{Bessel}s Elemente jedoch um 71 Minuten in der Länge von \textsc{Thulis}’ Beobachtungen ab. Das ist doch mehr, als man an Beobachtungsfehlern voraussetzen darf. Ich versuchte daher, den Fehler zu korrigieren, konnte aber dennoch einen Fehler von fast 29 Minuten in der Breite nicht vermeiden. Hier sind die Elemente

\begin{center}
\begingroup\fontsize{8}{10}\selectfont\setlength{\tabcolsep}{3pt}\renewcommand{\arraystretch}{1.1}
\begin{tabular}{lr}
Durchgang durch das Perihel im November . . . . & $17.746$ in Seeberg. \\
Länge des aufsteigenden Knotens . . . . . . . . & $340^\circ\ 11'$ \\
Neigung der Bahn . . . . . . . . . . . . . & $17\ 34$ \\
Länge der Sonnennähe . . . . . . . . . . . . & $157\ 17$ \\
Logarithmus des kleinsten Abstands . . . . . . . & $9.53969$ \\
Bewegung . . . . . . . . . . . . . . . . & rechtläufig \\
\end{tabular}\endgroup
\end{center}

„Diese Ergebnisse verdienen allerdings ebenfalls wenig Vertrauen.“
\SourceLeaf{271}
\par\indent Wir haben alle diese Umstände nur angeführt, um anderen, die sich an der Berechnung versuchen, vergebliche Mühe zu ersparen, falls sie die Beobachtungen aus Frankfurt und Marseille irgendwie in Übereinstimmung bringen wollen. Da \textsc{Bessel}s Bahn die Beobachtungen von \textsc{Olbers} bis zum 13. November so zufriedenstellend darstellt, dürfen diese Elemente als genau bestimmt gelten. Sollten wir künftig bessere Beobachtungen über einen längeren Zeitraum erhalten, dürften nur sehr geringe Verbesserungen nötig sein.

\hfill[\textsc{von Zach.}]

